\documentclass{article}
\usepackage{epsf, a4, makeidx, Manu}

\title{Index}

% Insertion d'un lien vers les sources
%
% \fichier{chemin}{nom}{ext} (eg \fichier{lib}{console}{c})
%
\newcommand{\fichier}[3]{
\href{../doxy/html/#2_8#3_source.html}{\lstinline!#1/#2.#3!}
(\href{../doxy/html/#2_8#3.html}{\sc dox})
}

\renewcommand{\labelitemi}{X}

% Pour le moment l'outil de traduction en md a du mal avec certains
% formats

\newcommand{\qemu}{{\em Qemu}}
\newcommand{\nasm}{{NASM} {}}
\newcommand{\iso}{{\sc iso}}
\newcommand{\manux}{{\em ManuX} {}}
\newcommand{\bios}{BIOS {}}
\newcommand{\pic}{PIC {}}
\newcommand{\irq}{IRQ {}}
\newcommand{\pc}{PC {}}

\author{Chaput Emmanuel}

\usepackage{graphicx} 
\usepackage{listings}
\input{lstasm}
\usepackage[french]{babel}
%\usepackage[utf8]{inputenc}


\begin{document}

\begin{center}
  \bf \em
   (Attention, la version actuelle de cette documentation contient des
liens cassés car l'outil mkdocs ne les gère pas correctement. Désolé,
je suis sur le coup !)
\end{center}

%===============================================================================
%
%===============================================================================
\section{ManuX-32}

   Bienvenue dans le monde merveilleux de ManuX ! 

%-------------------------------------------------------------------------------
%
%-------------------------------------------------------------------------------
\subsection{Allons-y, let's go, c'est parti les amis !}
   
   Le code source complet peut être téléchargé de la façon suivante~:

\begin{lstlisting}
 $ git clone git@github.com:ChaputEmmanuel/Manux.git
\end{lstlisting}

il sera alors compilé :

\begin{lstlisting}
 $ cd ManuX
 $ make
\end{lstlisting}

   puis on pourra le démarrer :

\begin{lstlisting}
 $ make run
\end{lstlisting}

   On peut également créer une image iso intégrant plusieurs noyaux
différents puis lancer le tout~:

\begin{lstlisting}
 $ make multiso
 $ make runiso
\end{lstlisting}

%-------------------------------------------------------------------------------
%
%-------------------------------------------------------------------------------
\subsection{Si ça marche, tant mieux, passe à la section suivante}

   La compilation et l’exécution du noyau nécessitent un certain
nombre d'outils que je vais tenter de lister ici brièvement et sur
lesquels je reviendrai dans la section \ref{section:outils} dédiée aux
outils.

   Pour être bref, sur une distribution Linux moderne, l'installation
des outils suivants devraient suffire~:
   
\begin{lstlisting}
 # apt install gcc xorriso nasm qemu-system-i386 mtools pandoc
\end{lstlisting}

%-------------------------------------------------------------------------------
%
%-------------------------------------------------------------------------------
\subsection{Ok, mais ensuite ?}

   Si vous avez réussi la première épreuve, vous avez donc un système
sur lequel vous pouvez compiler et faire tourner ManuX. Félicitation
et bienvenue dans mon monde \ldots

   Les différentes sections de ce document, forcément pas très à jour,
vous permettront de comprendre comment cela fonctionne et d'y apporter
vos modifications.

   Bon courage !
  

\end{document}

 
